% !TEX root = ../../cookbook.tex
\newpage
\recipe{Pasta alla Carbonara\label{carbonara}}{}{}{}
\vspace*{-0.5 cm}

\subsection*{Ingredienti \underline{per persona}}
\begin{ingredients}
	\item 80-120 g di pasta, secondo l' appetito prevedibile degli attavolati\footnote{Gli spaghetti sono la scelta tradizionale, ma io personalmente preferisco i rigatoni.}
	\item 2 tuorli, più un uovo intero ogni 2 persone\footnote{Diverse scuole di pensiero si scontrano riguardo l'uso di tuorli, uova intere, o una combinazione dei due. Si ottiene un'ottima Carbonara anche usando un uovo intero per persona, più un'uovo ``\textit{per la pentola}''. Se si favorisce la versione con i tuorli, gli albumi restanti possono essere usati per preparare delle ``\nameref{meringues}'' (p. \pageref{meringues}).}.
	\item 40 g di guanciale, tagliato a striscioline corte
	\item 15 g di Pecorino Romano grattuggiato
	\item Sale e Pepe
\end{ingredients}


\medskip

% --- Preparation ---
Cuocere la \textbf{pasta} poco più che al dente in acqua salata\footnote{Di regola, il livello di sale dell'acqua di cottura deve corrispondere a quello desiderato per la pasta che vi si cuoce.}.
Soffriggere il \textbf{guanciale} in una padella grande, che più tardi dovrà contenere tutta la pasta. 
Rimuovere il guanciale dalla padella quando dorato ma morbido, ma lasciarvi tutto il grasso e spegnere il fuoco.


Intanto, in una ciotola, mischiare insieme bene con una forchetta i \textbf{tuorli/uova}, il \textbf{Pecorino} grattuggiato, un poco di \textbf{sale}\footnote{Il pecorino e il guanciale provvederanno da soli un risultato naturalmente saporito.}, e \textbf{pepe} a volontà.

Una volta cotta la pasta, prelevare una ramaiolata (o un bicchierino) dell'acqua di cottura prima di scolare e versare la pasta nella padella con il grasso del guanciale. 

Versare un po' di acqua di cottura nella padella con la pasta e, mescolando continuamente, aggiungere il composto di uova e pecorino. 
Questa fase separerà una buona carbonara da una frittata: l'uovo dovrà addensarsi in'emulsione insieme al grasso e all'acqua, ed idealmente pastorizzarsi col calore residuo della padella, ma non dovrà assolutamente cuocere al punto di raggrumarsi\footnote{L'obbiettivo finale di una salsa cremosa si ottiene combinando abilmente una mescolatura continua, l'aggiunta di ulteriore acqua di cottura, e---solo se la pentola si è raffreddata al punto di non permettere alla salsa di addensarsi---la riaccensione del fuoco a fiamma bassissima.}. 

Ottenuta la consistenza desiderata, aggiungere il guanciale e servire con ulteriore pepe e pecorino grattuggiato.

%\vspace*{0.75 cm}

% --- Description ---
\begin{recipeblock}
	
\end{recipeblock}
